\subsection{Sự ra đời và tầm quan trọng của mạng LSTM}
Như đã phân tích ở chương trước, mạng RNN truyền thống gặp phải vấn đề nghiêm trọng về việc tiêu biến đạo hàm (Vanishing Gradient), khiến nó không thể ghi nhớ các thông tin trong một khoảng thời gian dài. Để khắc phục nhược điểm này, năm \textbf{1997}, Sepp Hochreiter và Jürgen Schmidhuber đã đề xuất kiến trúc \textbf{Long Short-Term Memory (LSTM)}. 

LSTM nhanh chóng trở thành một cuộc cách mạng trong lĩnh vực học sâu. Với khả năng lưu trữ thông tin qua hàng trăm bước thời gian, LSTM đã trở thành nền tảng cho các hệ thống dịch thuật tự động (Google Translate), nhận dạng giọng nói (Siri, Alexa) và đặc biệt là dự báo các biến động phức tạp trong kinh tế và kỹ thuật.

\subsection{Cấu trúc và cơ chế Cổng (Gates) của LSTM}
Khác với RNN chỉ có một cấu trúc đơn giản bên trong đơn vị lặp lại, LSTM sở hữu một kiến trúc tinh vi hơn nhiều với sự xuất hiện của \textbf{Cell State} và các \textbf{Cổng (Gates)}.

\subsubsection{Cell State (Trạng thái ô)}
Cell State có thể được hình dung như một "băng chuyền" thông tin chạy xuyên suốt toàn bộ chuỗi dữ liệu. Thông tin trên băng chuyền này rất ít bị thay đổi, giúp nó có thể truyền tải các thông tin quan trọng từ quá khứ xa xôi đến hiện tại mà không bị suy giảm đáng kể.

\subsubsection{Cơ chế các Cổng điều khiển}
LSTM sử dụng 3 loại cổng để điều chỉnh luồng thông tin ra vào Cell State:
\begin{itemize}
    \item \textbf{Forget Gate (Cổng quên):} Quyết định thông tin nào từ Cell State cũ sẽ bị xóa bỏ.
    \item \textbf{Input Gate (Cổng vào):} Quyết định thông tin mới nào sẽ được cập nhật vào Cell State.
    \item \textbf{Output Gate (Cổng ra):} Quyết định phần thông tin nào từ Cell State sẽ được đưa ra làm đầu ra và trạng thái ẩn mới.
\end{itemize}

\subsection{Công thức toán học của LSTM}
Tại mỗi bước thời gian $t$, quá trình tính toán trong một đơn vị LSTM được thực hiện qua các bước sau:

\begin{enumerate}
    \item \textbf{Cổng quên (Forget Gate):} Quyết định thông tin cũ cần xóa bỏ.
    \begin{equation}
        f_t = \sigma(W_f \cdot [h_{t-1}, x_t] + b_f)
    \end{equation}
    
    \item \textbf{Cổng vào (Input Gate):} Quyết định thông tin mới cần lưu trữ.
    \begin{equation}
        i_t = \sigma(W_i \cdot [h_{t-1}, x_t] + b_i)
    \end{equation}
    \begin{equation}
        \tilde{C}_t = \tanh(W_C \cdot [h_{t-1}, x_t] + b_C)
    \end{equation}
    
    \item \textbf{Cập nhật Trạng thái ô (Update Cell State):}
    \begin{equation}
        C_t = f_t * C_{t-1} + i_t * \tilde{C}_t
    \end{equation}
    
    \item \textbf{Cổng ra (Output Gate):} Xác định giá trị đầu ra.
    \begin{equation}
        o_t = \sigma(W_o \cdot [h_{t-1}, x_t] + b_o)
    \end{equation}
    \begin{equation}
        h_t = o_t * \tanh(C_t)
    \end{equation}
\end{enumerate}

Trong đó $\sigma$ là hàm \textit{sigmoid} điều tiết luồng thông tin trong khoảng [0, 1].

\subsection{Biến thể GRU (Gated Recurrent Unit) và So sánh với LSTM}
GRU là một biến thể hiện đại của LSTM, được giới thiệu bởi Cho và các cộng sự vào năm 2014. GRU đơn giản hóa cấu trúc của LSTM bằng cách kết hợp cổng quên và cổng vào thành một "cổng cập nhật" duy nhất, đồng thời bỏ đi trạng thái ô (Cell State).

\subsubsection{Sự khác biệt về cấu trúc:} 
Bảng dưới đây tóm tắt những điểm khác biệt cốt lõi giữa hai kiến trúc:

\begin{table}[H]
    \centering
    \begin{tabular}{|p{3cm}|p{5cm}|p{5cm}|}
        \hline
        \textbf{Đặc điểm} & \textbf{LSTM} & \textbf{GRU} \\ \hline
        \textbf{Số lượng cổng} & 3 cổng (Forget, Input, Output). & 2 cổng (Reset, Update). \\ \hline
        \textbf{Trạng thái} & Tách biệt Cell state ($C_t$) và Hidden state ($h_t$). & Chỉ sử dụng Hidden state ($h_t$). \\ \hline
        \textbf{Độ phức tạp} & Phức tạp hơn, nhiều tham số hơn. & Đơn giản hơn, ít tham số hơn. \\ \hline
        \textbf{Tốc độ huấn luyện} & Chậm hơn. & Nhanh hơn (thường nhanh hơn 20-30\%). \\ \hline
        \textbf{Khả năng học} & Tốt hơn trên các tập dữ liệu cực lớn và chuỗi rất dài. & Hiệu quả trên tập dữ liệu nhỏ/vừa hoặc chuỗi ngắn hơn. \\ \hline
    \end{tabular}
    \caption{So sánh kiến trúc LSTM và GRU}
\end{table}

\subsubsection{Tại sao GRU đôi khi hiệu quả hơn?}

Mặc dù LSTM về lý thuyết mạnh mẽ hơn, GRU thường được lựa chọn trong thực tế vì các lý do sau:
\begin{itemize}
    \item \textbf{Hiệu năng tính toán:} Với ít cổng hơn, GRU yêu cầu ít bộ nhớ và sức mạnh tính toán hơn, giúp huấn luyện nhanh hơn trên cùng một phần cứng.
    \item \textbf{Khả năng hội tụ:} Trên các tập dữ liệu không quá lớn, GRU có xu hướng hội tụ nhanh hơn và ít bị hiện tượng Overfitting hơn do mô hình đơn giản hơn.
    \item \textbf{Độ trễ thấp:} Trong các ứng dụng thời gian thực như chatbot hoặc nhận dạng giọng nói, tốc độ phản hồi của GRU là một lợi thế lớn.
\end{itemize}

\subsection{Dự báo chuỗi thời gian (Time Series Forecasting)}
Dự báo chuỗi thời gian là bài toán dự đoán các giá trị tương lai dựa trên các giá trị quan sát được trong quá khứ. Đây là lĩnh vực mà LSTM thể hiện sức mạnh vượt trội nhờ khả năng nắm bắt được cả xu hướng (trend) dài hạn và tính mùa vụ (seasonality).

\subsubsection{Kỹ thuật Cửa sổ trượt (Sliding Window)}
Để huấn luyện mô hình LSTM cho chuỗi thời gian, ta cần chuyển đổi dữ liệu về dạng học có giám sát (Supervised Learning). Kỹ thuật phổ biến nhất là sử dụng cửa sổ trượt: dùng $N$ giá trị trong quá khứ để dự đoán giá trị tại thời điểm tiếp theo.

\subsubsection{Tập dữ liệu Airline Passengers}
Báo cáo sử dụng tập dữ liệu kinh điển ghi lại số lượng hành khách hàng không hàng tháng từ năm 1949 đến 1960. Tập dữ liệu này có các đặc điểm: tính phi tuyến, tính xu hướng (Trend) và tính mùa vụ (Seasonality).

\subsection{Thực nghiệm với Python}
Dưới đây là mã nguồn xây dựng mạng LSTM để dự báo dữ liệu chuỗi thời gian lượng hành khách:

\subsubsection{Chuẩn bị dữ liệu}
\begin{lstlisting}[language=Python]
import numpy as np
import pandas as pd
import matplotlib.pyplot as plt
import math
from sklearn.preprocessing import MinMaxScaler
from sklearn.metrics import mean_squared_error
from tensorflow.keras.models import Sequential
from tensorflow.keras.layers import LSTM, Dense

print("--- 1. Loading Dataset ---")
file_path = "AirPassengers.csv"
df = pd.read_csv(file_path, usecols=[1])
data = df.values.astype('float32')

print("\n--- 2. Visualizing Raw Data ---")
plt.figure(figsize=(10, 6))
plt.plot(df, label='Monthly Passengers')
plt.title('International Airline Passengers')
plt.show()
\end{lstlisting}

\paragraph{Giải thích chi tiết:}
\begin{itemize}
    \item \textbf{\texttt{usecols=[1]}:} Tập dữ liệu gốc chứa cột thời gian và cột số lượng hành khách. Vì LSTM sẽ học dựa trên vị trí thứ tự của mảng, ta chỉ cần cột giá trị.
    \item \textbf{Trực quan hóa dữ liệu thô:} Bước này giúp xác định xu hướng và mùa vụ của dữ liệu trước khi chọn kiến trúc mạng.
\end{itemize}

\begin{figure}[H]
    \centering
    \includegraphics[width=1\linewidth]{images/41.png}
    \caption{Dữ liệu gốc lượng hành khách hàng tháng}
    \label{fig:placeholder}
\end{figure}

Biểu đồ đường (Line chart) trực quan hóa sự biến động của số lượng hành khách hàng không quốc tế hàng tháng. Dựa vào các thông tin trên biểu đồ, ta có thể rút ra các đặc điểm sau:

\begin{itemize}
    \item \textbf{Thông tin các trục:} Trục hoành \textit{Month (Index)} thể hiện các bước thời gian từ $0$ đến hơn $140$. Trục tung \textit{Passengers} thể hiện số lượng hành khách tương ứng.
    \item \textbf{Xu hướng tăng (Upward Trend):} Đường biểu diễn \textit{Monthly Passengers} (màu xanh) cho thấy sự tăng trưởng liên tục và rõ rệt. Dữ liệu bắt đầu ở mức xấp xỉ $100$ hành khách (tại index $0$) và đạt đỉnh cao nhất vượt mức $600$ (gần index $140$).
    \item \textbf{Tính mùa vụ (Seasonality):} Dữ liệu liên tục dao động tạo thành các đỉnh (cao điểm) và đáy (thấp điểm) lặp lại đều đặn dọc theo trục thời gian.
    \item \textbf{Biên độ dao động mở rộng (Increasing Variance):} Khoảng cách giữa các đỉnh và đáy trong cùng một chu kỳ ngày càng giãn rộng khi tiến về phía cuối biểu đồ. Điều này chỉ ra rằng đây là một chuỗi thời gian dạng nhân (Multiplicative Time Series) chứ không phải dạng cộng.
\end{itemize}

\begin{lstlisting}[language=Python]
print("\n--- 3. Preprocessing ---")
scaler = MinMaxScaler(feature_range=(0, 1))
dataset = scaler.fit_transform(data)

train_size = int(len(dataset) * 0.7)
test_size = len(dataset) - train_size
train, test = dataset[0:train_size,:], dataset[train_size:len(dataset),:]

def create_dataset(dataset, look_back=1):
    dataX, dataY = [], []
    for i in range(len(dataset)-look_back-1):
        a = dataset[i:(i+look_back), 0]
        dataX.append(a)
        dataY.append(dataset[i + look_back, 0])
    return np.array(dataX), np.array(dataY)

look_back = 12
trainX, trainY = create_dataset(train, look_back)
testX, testY = create_dataset(test, look_back)

trainX = np.reshape(trainX, (trainX.shape[0], 1, trainX.shape[1]))
testX = np.reshape(testX, (testX.shape[0], 1, testX.shape[1]))
\end{lstlisting}

\paragraph{Giải thích chi tiết:}
\begin{itemize}
    \item \textbf{\texttt{MinMaxScaler(0, 1)}:} Đưa dữ liệu về khoảng [0, 1] để phù hợp với hàm kích hoạt của LSTM.
    \item \textbf{Hàm \texttt{create\_dataset} (Sliding Window):} Chuyển đổi chuỗi thành bài toán học có giám sát.
    \item \textbf{Tham số \texttt{look\_back}:} Xác định số lượng bước thời gian trong quá khứ được dùng để dự báo. Ở đây ta chọn 12 tháng gần nhất để dự đoán tháng tiếp theo.
\end{itemize}

\subsubsection{Xây dựng và Huấn luyện mô hình LSTM}
\begin{lstlisting}[language=Python]
print("\n--- 4. Building LSTM Model ---")
model = Sequential([
    LSTM(50, return_sequences=True, input_shape=(1, look_back)),
    Dropout(0.5),
    LSTM(50),
    Dropout(0.5),
    Dense(1)
])
model.compile(loss='mean_squared_error', optimizer='adam')
model.summary()

print("\n--- 5. Training Model ---")
history = model.fit(trainX, trainY, epochs=100, batch_size=1, verbose=0, validation_split=0.1)
\end{lstlisting}

\paragraph{Giải thích cấu trúc và tham số:}
\begin{itemize}
    \item \textbf{\texttt{return\_sequences=True}:} Truyền toàn bộ chuỗi trạng thái ẩn từ lớp LSTM thứ nhất sang lớp LSTM thứ hai để trích xuất đặc trưng sâu hơn.
    \item \textbf{\texttt{Dropout(0.5)}:} Tắt ngẫu nhiên $50\%$ nơ-ron trong quá trình huấn luyện nhằm ngăn chặn triệt để hiện tượng Overfitting.
    \item \textbf{Hàm mất mát \& Tối ưu hóa:} Sử dụng \texttt{mean\_squared\_error} phù hợp cho bài toán hồi quy (Regression), kết hợp bộ tối ưu \texttt{adam} để tăng tốc độ hội tụ.
    \item \textbf{\texttt{batch\_size=1}:} Cập nhật trọng số ngay sau từng mẫu dữ liệu (Online Learning), giúp mô hình bám sát các biến động chi tiết của chuỗi thời gian.
\end{itemize}

\begin{figure}[H]
    \centering
    \includegraphics[width=1\linewidth]{images/42.png}
    \caption{Tham số mô hình}
    \label{fig:placeholder}
\end{figure}

\begin{figure}[H]
    \centering
    \includegraphics[width=1\linewidth]{images/43.png}
    \caption{Đồ thị Loss trong quá trình huấn luyện LSTM}
    \label{fig:placeholder}
\end{figure}

Biểu đồ \textit{LSTM Training History} thể hiện sự biến thiên của hàm mất mát (Loss) qua 150 vòng huấn luyện (epochs):

\begin{itemize}
    \item \textbf{Tập huấn luyện (Train Loss):} Giảm mạnh ở những epoch đầu và nhanh chóng hội tụ, duy trì ổn định ở mức cực kỳ thấp (xấp xỉ 0.001 - 0.002).
    
    \item \textbf{Tập kiểm định (Val Loss):} Điểm khác biệt lớn nhất là đường Val Loss giờ đây \textbf{bám rất sát} đường Train Loss. Mặc dù vẫn còn vài dao động nhỏ, biên độ các đỉnh đã giảm đáng kể và không còn hiện tượng lệch pha kéo dài.
    
    \item \textbf{Kết luận:} Việc áp dụng kỹ thuật Dropout đã khắc phục thành công hiện tượng Overfitting. Khoảng cách (gap) giữa hai tập dữ liệu gần như bị xóa bỏ, chứng tỏ mô hình đã đạt được khả năng tổng quát hóa rất tốt trên dữ liệu chưa từng gặp.
\end{itemize}

\subsubsection{Đánh giá kết quả Dự báo}
\begin{lstlisting}[language=Python]
print("\n--- 6. Predicting & Evaluating ---")
trainPredict = model.predict(trainX)
testPredict = model.predict(testX)

# Inverse transforms
trainPredict = scaler.inverse_transform(trainPredict)
trainY_orig = scaler.inverse_transform([trainY])
testPredict = scaler.inverse_transform(testPredict)
testY_orig = scaler.inverse_transform([testY])

# Calculate RMSE
trainScore = math.sqrt(mean_squared_error(trainY_orig[0], trainPredict[:,0]))
print(f'Train Score: {trainScore:.2f} RMSE')
testScore = math.sqrt(mean_squared_error(testY_orig[0], testPredict[:,0]))
print(f'Test Score: {testScore:.2f} RMSE')
\end{lstlisting}

\paragraph{Giải thích chi tiết:}
\begin{itemize}
    \item \textbf{\texttt{inverse\_transform}:} Đưa kết quả dự báo quay lại thang đo gốc (số lượng hành khách).
    \item \textbf{Chỉ số RMSE:} Căn bậc hai của sai số bình phương trung bình, giúp ta biết sai số trung bình là bao nhiêu hành khách. Đây là chỉ số quan trọng nhất trong dự báo chuỗi thời gian, đơn vị của nó khớp với đơn vị gốc của dữ liệu.
\end{itemize}

\begin{figure}[H]
    \centering
    \includegraphics[width=1\linewidth]{images/44.png}
    \caption{Kết quả dự báo lượng hành khách trên tập Train và Test với look\_back = 12}
    \label{fig:placeholder}
\end{figure}

\subsection{Phân tích kết quả thực nghiệm và Đánh giá dự báo}

Kết quả dự báo lượng hành khách hàng không bằng mô hình LSTM cho thấy khả năng nắm bắt quy luật thời gian ấn tượng của mạng nơ-ron hồi tiếp cải tiến này.

\subsubsection{Phân tích chỉ số RMSE (Root Mean Squared Error)}
Chỉ số RMSE là thước đo tiêu chuẩn trong các bài toán hồi quy chuỗi thời gian:
\begin{itemize}
    \item \textbf{Ý nghĩa:} RMSE đo lường độ lệch trung bình giữa giá trị hành khách dự báo và giá trị hành khách thực tế. Với dữ liệu hàng không, một giá trị RMSE thấp cho thấy mô hình có sai số rất nhỏ so với biên độ tổng lượng hành khách hàng tháng.
    \item \textbf{Tính tổng quát hóa:} Việc RMSE trên tập Train và tập Test duy trì ở mức ổn định, không quá chênh lệch chứng tỏ mô hình LSTM đã học được "quy luật" cốt lõi của chuỗi dữ liệu chứ không rơi vào tình trạng ghi nhớ máy móc (Overfitting).
\end{itemize}

\subsubsection{Đánh giá trực quan qua biểu đồ dự báo}
Biểu đồ so sánh trực tiếp dữ liệu thực tế (Actual Data - màu xanh dương) với kết quả dự báo trên tập huấn luyện (Train Prediction - màu cam) và tập kiểm tra (Test Prediction - màu xanh lá):
\begin{itemize}
    \item \textbf{Khớp dữ liệu quá khứ (Train):} Đường màu cam bám rất sát với đường thực tế. Mô hình đã nhận diện xuất sắc cả xu hướng tăng trưởng tuyến tính (Trend) và chu kỳ dao động lặp lại (Seasonality).
    \item \textbf{Dự báo tương lai (Test):} Trên tập dữ liệu chưa từng gặp (phần đường màu xanh lá ở cuối đồ thị), mô hình vẫn bám sát hoàn hảo các biến động. Đặc biệt, LSTM đã dự báo vô cùng chính xác các đỉnh cực đại (vượt ngưỡng 600 hành khách) và đáy của chu kỳ mở rộng.
    \item \textbf{Kết luận:} Mô hình không hề bị trễ pha (lagging) - một lỗi rất phổ biến khi dùng mạng Neural dự báo chuỗi thời gian. Điều này khẳng định cấu trúc LSTM kết hợp Dropout đã mang lại khả năng dự báo ngoại suy xuất sắc cho dữ liệu phi tuyến tính.
\end{itemize}

\subsection{Thử nghiệm tham số: Tầm quan trọng của look\_back}

Tham số \texttt{look\_back} xác định số lượng bước thời gian trong quá khứ được dùng làm ngữ cảnh để dự báo tương lai. Thực nghiệm tiến hành đánh giá ba cấu hình:

\begin{itemize}
    \item \textbf{Cấu hình 1 (look\_back = 1):} Chỉ sử dụng dữ liệu 1 tháng liền trước.
    \item \textbf{Cấu hình 2 (look\_back = 3):} Sử dụng chuỗi dữ liệu 3 tháng gần nhất (1 quý).
    \item \textbf{Cấu hình 3 (look\_back = 12):} Sử dụng toàn bộ chuỗi dữ liệu 12 tháng gần nhất (1 năm).
\end{itemize}

\subsubsection{Bảng so sánh sai số RMSE:}

\begin{table}[H]
    \centering
    \begin{tabular}{|l|c|c|}
        \hline
        \textbf{Cấu hình} & \textbf{Train RMSE} & \textbf{Test RMSE} \\ \hline
        \textbf{look\_back = 1} & 24.21 & 65.41 \\ \hline
        \textbf{look\_back = 3} & 22.16 & 55.22 \\ \hline
        \textbf{look\_back = 12} & 12.62 & 22.06 \\ \hline
    \end{tabular}
    \caption{So sánh sai số dự báo theo tham số look\_back}
\end{table}

\subsubsection{Phân tích:} 
Việc mở rộng \texttt{look\_back} từ $1$ lên $3$ giúp mạng LSTM nắm bắt tốt hơn "đà" (momentum) biến động của dữ liệu. Đặc biệt, khi tăng \texttt{look\_back = 12}, sai số giảm mang tính đột phá trên cả hai tập. Điều này hoàn toàn phù hợp với đặc thù dữ liệu hàng không: tính mùa vụ lặp lại theo chu kỳ đúng 12 tháng. Cửa sổ thời gian 1 năm giúp mô hình "nhìn" thấy được toàn cảnh quy luật của năm trước (tháng cao điểm, tháng thấp điểm) để ánh xạ và dự báo cực kỳ chính xác cho năm tiếp theo.

\subsection{Kết luận chung}

Nhờ cơ chế cổng nhớ thông minh (memory gates) và việc thiết lập tham số cửa sổ thời gian (\texttt{look\_back}) bám sát chu kỳ thực tế, mô hình LSTM đã chứng minh năng lực xuất sắc. Nó không chỉ duy trì được bộ nhớ về xu hướng dài hạn (Long-term Trend) mà còn thích ứng hoàn hảo với tính mùa vụ (Seasonality) của ngành hàng không. Mức sai số RMSE thấp trên tập kiểm thử khẳng định LSTM là một phương pháp tiếp cận mạnh mẽ và tối ưu cho bài toán dự báo chuỗi thời gian phi tuyến tính.